\documentclass{article}
\usepackage[T1]{fontenc}
\usepackage{lmodern}
\usepackage{graphicx}
\usepackage{amsmath,amssymb}
\usepackage[a4paper,margin=2cm]{geometry}
\usepackage[absolute,overlay]{textpos}
\setlength{\TPHorizModule}{1mm}
\setlength{\TPVertModule}{1mm}
\usepackage[indonesian]{babel}
\usepackage{hyperref}
\setlength{\emergencystretch}{2em}
% unit: Halaman 14

\begin{document}

% === Halaman 14 (python/native) ===

Rinaldi Munir/II4021 Kriptografi

14

Medan Galois (Galois Field)

• Medan Galois adalah medan berhingga dengan pm elemen, p adalah bilangan prima 

dan m  1.

• Notasi: GF(pm)

• Kasus paling sederhana: bila m = 1 → GF(p) atau Fp 

• GF(p) terdiri dari himpunan \{0, 1, 2, …, p – 1\} dan operasi penjumlahan dan 

perkalian dilakukan di dalam modulus p. 

   1. Penjumlahan, dilakukan dalam modulus p

     Jika a, b  GF(p), maka a + b = r, yang dalam hal ini  r = (a + b) mod p,  0  r  p – 1

 2. Perkalian, dilakukan dalam modulus p

    Jika a, b  GF(p), maka a  b = s, yang dalam hal ini s = (a  b) mod p,   0  s  p – 1

\end{document}
